\section{Materials and Methods}

\subsection{Task Formulation and Scope}\label{sec:problem}

We study \emph{rule-governed, evidence-bearing document review}. An instance consists of a source-addressable document $d$, an explicit rule package $K$, a policy $P$ for aggregating findings, and an accountable reviewer. The system prepares rule-linked findings and unresolved items; the reviewer determines their disposition and the final opinion. This task is broader than technical proposal review but narrower than arbitrary document understanding: it requires explicit criteria, identifiable evidence sources and a meaningful distinction between a machine recommendation and an accountable decision. A source range identifies what was read, not whether its contents are true or sufficient.

The empirical instance is preliminary review of Chinese technical solution proposals before construction. Its primary business object is the \emph{adequacy of the proposed design}, not proof that the proposed implementation has already been deployed. A rule must specify the lifecycle stage, applicability, expected evidence and predicate being checked. We retain the released package and labels for the retrospective evaluation, rather than silently relabel them to fit this formulation. Where its predicates conflate design and implementation, or parser support with business compliance, the discrepancy is reported and must be resolved before a revised package is evaluated.

Four requirements guide the design: rule-linked evidence and explicit non-detection records; rule-specific eligibility to contribute a finding; a visible distinction between confirmed findings, unresolved items and final disposition; and inspectable feedback with controlled activation and recoverable configuration. These are design requirements, not four completed implementation guarantees. Table~\ref{tab:reqs} separates the evaluated prototype from the stronger deployment contract.

\begin{table}[!ht]
\caption{Requirements and evidence status. The evaluated prototype does not yet implement the complete approval and recovery contract.}
\label{tab:reqs}
\centering\footnotesize
\begin{tabularx}{\textwidth}{@{}p{2.4cm}XX@{}}
\toprule
Requirement & Evaluated prototype & Remaining deployment requirement \\
\midrule
Evidence traceability & Rule identifiers, source ranges, verifier absence records and recorded gate outcomes & Separate source location, scope relation and semantic support; validate converted inputs \\
Rule-specific eligibility & Hard/hybrid/soft routing and package validation & Validate effective authority after every update, not only the tier field \\
Accountable disposition & Confirmed/pending states, deterministic recommendation and recorded decisions & Host-authenticated approval; explicit disposition or waiver of pending items before signing \\
Feedback and recovery & Overrides, statistics, exemplars, audit records and an override-count version string & Content-addressed immutable configurations, approval binding and demonstrated restoration \\
\bottomrule
\end{tabularx}
\end{table}

\subsection{Operational Authority Types}\label{sec:rules}

Each YAML rule declares an identifier, dimension, severity, tier, question, rationale, remediation, applicability condition, verifier specification and status. A shared lexicon supplies heading aliases, placeholder and negation tokens, and keyword families. The package contains 28 rules: 10 hard, 13 hybrid and 5 soft. Seven have applicability preconditions; the lexicon contains 25 keyword families, 11 heading groups, 12 placeholder tokens and 11 negation tokens.

The tier is an \emph{operational routing declaration}. A hard rule uses a decisive implemented predicate; a hybrid rule combines deterministic checks or candidate localisation with a judgement step; a soft rule raises an unresolved item. Decisive execution of a predicate does not prove that the predicate fully captures the business requirement. For example, a lexical declaration of a load-test plan does not establish its scope, responsibilities, timing or performance criteria. Assigning a hard tier to that proxy cannot make those unmeasured properties automatically verifiable. Domain owners must first define the requirement and then specify which subpredicates the verifier actually covers.

\begin{table}[!ht]
\caption{Default routes in the evaluated package. These routes describe finding eligibility under the frozen configuration, not authentication of human approval.}
\label{tab:tiers}
\centering\small
\begin{tabularx}{\textwidth}{@{}lXXX@{}}
\toprule
Tier & Verifier role & Default outcomes & Finding route \\
\midrule
hard & Decisive configured predicate & pass, fail, not applicable & Verifier-generated finding; model registration rejected \\
 hybrid & Structural checks and candidate localisation & structural result, candidate, pending, not applicable & Structural failure may be automatic; otherwise a gate-accepted judgement may count \\
soft & No verifier & unresolved item & Model registration remains pending until an explicit disposition \\
\bottomrule
\end{tabularx}
\end{table}

\begin{figure}[!ht]
\centering
\begin{tikzpicture}[font=\small,>=Stealth,node distance=0.5cm,
box/.style={draw,rounded corners,align=center,text width=3.7cm,minimum height=1.5cm}]
\node[box] (h) {\textbf{Hard}\configured predicate\verifier finding};
\node[box,right=of h] (y) {\textbf{Hybrid}\structural check / candidate\gate-accepted judgement};
\node[box,right=of y] (s) {\textbf{Soft}\unresolved item\reviewer disposition};
\node[box,text width=8cm,below=of y] (p) {Recommendation aggregates eligible findings;\final signing needs a separate authenticated approval contract};
\draw[->] (h.south) |- (p.west);\draw[->] (y)--(p);\draw[->,dashed] (s.south) |- (p.east);
\end{tikzpicture}
\caption{Operational routes and their boundary. A tier is one component of effective authority, not a sufficient permission ceiling for subsequent edits. The complete signing contract is a deployment requirement, not an evaluated capability.}
\label{fig:type}
\end{figure}

In the default registration path, absence evidence is verifier-owned: a model cannot establish non-detection by inventing a quotation such as ``not mentioned''. A verifier absence record identifies the expected section and the failed search predicate; it does not establish semantic absence outside the parser's supported input and vocabulary. Hybrid non-detection defaults to pending. Candidate existence is not satisfaction, and lack of a registered finding is not a verified pass.

Package validation rejects inconsistent verifier/tier combinations, missing lexicon references and duplicate identifiers. This is useful static validation but does not enforce the same restrictions after every learned edit. The evaluated update paths can change no-candidate behaviour, activate a valid hybrid draft, alter severity and reset overrides without preserving a complete effective-authority ceiling. Consequently, immutability of the tier field alone does not establish bounded authority.

All four blocker rules in this package (CAP-001, SEC-001, SEC-002 and SEC-004) are hybrid; the ten hard rules are major or minor. This distribution motivates tier-specific failure analysis, but is a characteristic of this checklist, not a general inverse relationship between severity and verifiability.

\begin{table}[!ht]
\caption{Composition of the frozen package used in the retrospective evaluation.}
\label{tab:composition}
\centering\footnotesize
\begin{tabular}{@{}lrrr rrrr@{}}
\toprule
Dimension & Hard & Hybrid & Soft & Total & Blocker & Major & Minor \\
\midrule
Capacity & 2 & 3 & 1 & 6 & 1 & 5 & 0 \\
Security & 0 & 4 & 1 & 5 & 3 & 2 & 0 \\
Interface & 2 & 3 & 0 & 5 & 0 & 3 & 2 \\
Documentation & 6 & 0 & 2 & 8 & 0 & 4 & 4 \\
Channel & 0 & 1 & 0 & 1 & 0 & 1 & 0 \\
Data synchronisation & 0 & 1 & 0 & 1 & 0 & 1 & 0 \\
Configuration & 0 & 0 & 1 & 1 & 0 & 0 & 1 \\
Monitoring & 0 & 1 & 0 & 1 & 0 & 0 & 1 \\
\midrule
Total & 10 & 13 & 5 & 28 & 4 & 16 & 8 \\
\bottomrule
\end{tabular}
\end{table}

\subsection{Parsing, Verifiers and Business Predicates}\label{sec:verifiers}

The review engine accepts Markdown. Word and PDF conversion occurs upstream through the host's office tooling; no conversion-fidelity evaluation is reported here. The parser constructs a heading tree, tables and Mermaid sequence diagrams~\cite{mermaid}, with source offsets. Its normalisation includes NFKC folding, punctuation folding, emphasis removal and whitespace collapsing, with a map back to the source. Chinese and Latin sentence terminators and table-cell boundaries define candidate units. Absence records use an expected section and a zero-length source range.

\begin{table}[!ht]
\caption{Implemented verifier families. An established parser predicate is not necessarily the complete business predicate. Uses count occurrences in verifier trees.}
\label{tab:verifiers}
\centering\footnotesize
\begin{tabularx}{\textwidth}{@{}p{2.4cm}rXX@{}}
\toprule
Family & Uses & Implemented check & Scope limitation \\
\midrule
statement & 16 & Keyword matches, local negation, placeholders, assertion cues & Does not establish semantic sufficiency \\
section & 9 & Heading, body length and placeholder checks & Alias and format dependent \\
table & 3 & Column groups, row count and placeholder cells & Not correctness of cell content \\
mermaid-sequence & 3 & Parseable diagram, participants, messages and required families & Cannot inspect arbitrary embedded images \\
cross-ref & 3 & Attachment-reference resolution & Not attachment correctness \\
quantity & 1 & Number/unit near a context family & Not validity of the quantity \\
all, first & 4, 2 & Conjunction and ordered locating fallback & Inherits child limitations \\
\bottomrule
\end{tabularx}
\end{table}

Negation checks use a window around a matched family; placeholder checks combine lexical tokens and structural cues. These choices recover defects that substring presence misses, but their correctness remains relative to the declared predicate and lexicon. The frozen diagram rules report an image-only diagram as a failure because they require Mermaid. That result supports compliance with a \emph{declared machine-readable format}, not a finding that the depicted design is missing or wrong. Unless an organisation requires that format, unsupported image input should become \emph{needs manual review}, rather than business noncompliance. This correction changes rule behaviour and may require label migration and reevaluation; the retrospective counts below are not presented as results of such a corrected package.

\subsection{Located-Evidence Gate}\label{sec:gate}

Model findings use a registration tool. The gate rejects hard-rule registration, disallows model-produced absence evidence, checks normalised quote length (8--300 characters), locates the quotation and checks its relation to a hybrid candidate. Unknown or soft-rule submissions remain pending rather than gaining automatic eligibility. Rejection reasons are stored for inspection and possible repair.

\begin{algo}[H]
\caption{Gate in the evaluated prototype (simplified). Scope membership is overlap, not complete containment; location does not establish entailment.}
\label{alg:gate}
\begin{algorithmic}[1]
\Require document, rule id, evidence list, candidate ranges, optional stated section
\Ensure gate-confirmed finding or pending item with reasons
\State Resolve rule; unknown or soft rule adds a pending reason
\If{tier is hard} \State Reject model registration \EndIf
\If{evidence is empty} \State Add missing-evidence reason \EndIf
\For{each evidence item $e$}
  \If{$e$ asserts absence or starts with an absence marker}
    \State Add absence-not-allowed reason
  \EndIf
  \State Normalise $e.\text{quote}$; add a reason if length is outside 8--300
  \State Find all source hits of the normalised quote
  \State Prefer candidate-and-section hit, then candidate hit, then section hit, then first hit
  \If{no hit exists} \State Add quote-not-found reason
  \Else
    \State Attach the selected source range $h$
    \If{tier is hybrid and $h$ overlaps no candidate} \State Add out-of-candidates reason \EndIf
  \EndIf
\EndFor
\State \Return gate-confirmed if no reason; otherwise pending with reasons
\end{algorithmic}
\end{algo}

The gate establishes recoverable quotation and a positional relationship, not that the quotation supports the allegation. In particular, partial overlap can pass, and a genuine but misleading quotation can be accepted. A strict containment contract would instead require every quoted character to lie within an allowed range and would need separate boundary tests and a new evaluation. We do not attribute zero false positives or prompt-injection resistance to this gate. The reported 42/53 fraction is \emph{gate acceptance}, not semantic evidence validity.

\subsection{Recommendation and Final Disposition}\label{sec:policy}

Let $C$ contain findings with an eligible route (hard, hybrid or expert) and status confirmed, accepted or rewritten. Under the frozen policy, gate-confirmed hybrid findings can already enter $C$ before expert adoption. With $m$ the number of minor findings in $C$ and $\theta=3$,
\begin{equation*}
P(C)=\begin{cases}
\text{Fail},&\exists i\in C:\operatorname{severity}(i)=\text{blocker},\\
\text{Revise},&\exists i\in C:\operatorname{severity}(i)=\text{major}\ \text{or}\ m\ge\theta,\\
\text{Pass},&\text{otherwise}.
\end{cases}
\end{equation*}
Pending and soft items do not count; dismissed findings leave $C$. Opinion-channel findings become numbered clauses; checklist-channel findings form a correction annex, but both eligible channels contribute to the policy. Deterministic aggregation means repeatability for the same finding set; it does not make that set correct, and a model can influence the recommendation through accepted registrations or severity choices.

The prototype records a final verdict through a tool requiring a reviewer-name field. That field is provenance metadata, not authentication: an agent can supply a name. Finalisation also does not reliably require every pending item to have been disposed of. The current implementation therefore supports a \emph{trusted-operator workflow}, not a demonstrated software-enforced human-signature boundary. Suggested/final verdict disagreement may be recorded, but cannot by itself measure rubber-stamping.

\subsection{Feedback Records and Configuration Recovery}\label{sec:memory}\label{sec:feedback}

The review preset prescribes preliminary review, expert disposition and finalisation. It asks the agent to present findings in small batches with quotes and suggested remedies; accept, dismiss, rewrite and add actions become structured records. Statistics and exemplars can subsequently be recalled. The software stores feedback; this is not evidence that later judgements improve, or that the person recorded as operator genuinely approved each action.

\begin{table}[!ht]
\caption{Storage objects in the prototype. Write paths describe implemented operations, not independently authenticated human acts.}
\label{tab:storage}
\centering\footnotesize
\begin{tabularx}{\textwidth}{@{}p{3cm}XX@{}}
\toprule
Object & Contents and write operation & Subsequent use \\
\midrule
rule\_overrides & Status, severity, suggestions, whitelisted parameters, drafts; update tools & Effective rule loading and rules view \\
reviews & Input digest, recommendations, final value, timings, displayed version & Resumption and dashboard \\
issues & Findings, evidence, tier, severity and gate outcome; check/registration/decision & Opinion and evaluation \\
feedback & Accept/dismiss/rewrite records with reasons & Aggregation \\
exemplars & Accepted suggestions and rewrites; aggregation & Suggestion context \\
rule\_candidates & Unmatched additions and draft rules & Rule inspection/activation \\
eval\_runs & Configuration results and per-rule outputs & Dashboard and retrospective analysis \\
\bottomrule
\end{tabularx}
\end{table}

A requested severity or status edit can carry a verbatim \texttt{basis.quote}, recorded with operator, time and review. An arbitrary sufficiently long string satisfies this check; it does not prove consent. Parameter changes and valid-draft activation can also change decision eligibility while leaving the tier field unchanged. Reset removes overrides and keeps statistics, but may restore stronger shipped settings than the last approved settings.

The displayed identifier \texttt{pack@ver+ovN} counts rules with active overrides. Two distinct configurations can have the same $N$; it is not a content identity and does not enable exact historical replay. Audit records are useful provenance, but reversibility requires retaining the actual configuration to be restored and validating its authority, rather than merely deleting overrides.

\subsection{Required Human-Approval and Recovery Contract}\label{sec:contract}

For deployment under a software-enforced human-approval claim, we specify a stronger contract, not an already implemented extension. Effective authority is a tuple of permitted decision subject, admissible evidence, no-candidate behaviour, severity, applicability, activation status and policy eligibility. Every write path must validate the merged configuration against its approved ceiling, including parameter edits, draft activation, finding severity, enable/disable and reset.

A trusted host must issue an approval event after displaying the exact object and action to an authenticated reviewer. The event must bind review/rule identity, input and opinion/configuration content, action and values, principal and consumption or expiry state. An agent-provided name or quotation must not mint that event. Content changes, cross-object use and replay invalidate approval. Final signing requires explicit disposition of pending items or a separately authorised waiver with a reason; the interface must not confuse ``Pass on confirmed findings'' with review completion.

Recoverable versions require content-addressed immutable snapshots of rules, lexicon, overrides, policy and relevant parser/verifier versions, linked to input and approval records. Restoration must recover the specified approved state and reproduce its expected behaviour on fixed inputs while retaining the audit history. Acceptance tests must include valid approvals as well as forged, missing, replayed and changed-content approvals; authority-changing parameter/draft/severity/reset paths; unresolved-item signing; partial-overlap quotes; and distinct configurations with the same override count. These tests and the resulting repairs remain future engineering work in this manuscript.

This contract assumes trusted host authentication, storage and the controlled review APIs. A malicious administrator or an agent with arbitrary host shell/file write access is outside the plugin boundary. Scoped reading may reduce exposed context, but does not establish protection against prompt injection, semantic manipulation or exfiltration~\cite{shi2024inject,mempriv2025}.

\subsection{Implementation and Evaluated Version}\label{sec:impl}

HARK~\cite{artifact} is implemented as a plugin for a conversational agent harness supplying model access, file conversion, structured questions, delivery, JSON-backed storage and a web server. Its pure-function engine contains the parser, verifiers, gate, policy, aggregation and renderer. The plugin has seven storage objects, nine tools, three state-query slash commands and a host-served dashboard; it needs no separate application backend or embedding index. Host access control is reused, but operation-specific approval is a separate requirement and is not supplied by a reviewer-name field.

\begin{table}[!ht]
\caption{Principal tool routes and their current boundaries. A prompt instruction is not an authentication mechanism.}
\label{tab:tools}
\centering\footnotesize
\begin{tabularx}{\textwidth}{@{}p{3.4cm}XX@{}}
\toprule
Tool & Purpose & Evaluated boundary \\
\midrule
\texttt{review\_parse} & Parse/resume input & Supported text and input digest \\
\texttt{review\_check} & Run configured verifiers & Package validation; hard findings/hybrid candidates \\
\texttt{review\_record\_issue} & Register a judgement & Hard registration rejected; soft/invalid evidence pending; hybrid overlap gate \\
\texttt{review\_status} & Show findings and unresolved items & Visibility does not force disposition \\
\texttt{review\_decide} & Record disposition/final value & Name is recorded, not authenticated consent \\
\texttt{review\_finalize} & Render, close and aggregate & Pending completion is not a reliable enforced prerequisite \\
\texttt{rules\_list}, \texttt{rules\_propose}, \texttt{memory\_recall} & Inspect, edit/draft and recall & Tier-field restrictions do not cover all effective-authority changes \\
\bottomrule
\end{tabularx}
\end{table}

The retrospective analysis uses released outputs associated with repository snapshot \texttt{b566b08de2976451747215af4dc874a0846f6173}. It does not mix those results with the proposed approval/recovery contract. The uploaded manuscript differs from the repository manuscript; the revised paper therefore identifies the code snapshot explicitly rather than treating the texts as identical. The stored outputs were reaggregated without changing inputs, labels, rules or runtime behaviour. A future rule, gate, policy or state-machine repair must be evaluated under a new frozen version; these historical detection scores do not certify it.

The artifact supplies tests, but a clean-environment full-suite reproduction is not established here: the checked environment encountered missing host-peer dependencies and a packaging-ignore assertion failure. We therefore do not assert an all-passing published test count. Moving to another document domain requires new predicates, source adapters and labelled evaluation, not merely a second YAML file.

